\section{Introduction}
\label{sec:introduction}

The \gls{NSF}--\gls{DOE} Vera C. Rubin Observatory is a ground-based, wide-field optical-near-infrared facility on Cerro Pachón in northern Chile, built to carry out the \gls{LSST} \citep[][]{2019ApJ...873..111I}.
The observatory is named in honor of Vera C. Rubin, whose observations in the latter half of the 20th century established the first compelling empirical case for dark matter \citep[][]{1970ApJ...159..379R, 1980ApJ...238..471R}.
Its 10-yr survey will image the entire visible southern sky every 3--4 nights, reaching a coadded depth of $\sim$\,~27\,mag in the $r$-band across an 18,000\,deg$^{\rm 2}$ footprint \citep[][]{2022ApJS..258....1B}.

At the heart of the facility is the \lsstcampixels-gigapixel LSSTCam, mounted on the \primarymirrordiameter (\effectiveaperture effective aperture) Simonyi Survey Telescope \citep[][]{2024SPIE13094E..09S} .
Together they deliver an instantaneous field of view of \rubinfov in six broadband filters, $ugrizy$ (320--1050\,nm), with seeing-limited image quality.
The system's \'{e}tendue of \rubinetendue exceeds that of any previous optical facility by over an order of magnitude, enabling the rapid, deep, all-sky cadence that underpins Rubin's four core science themes: the nature of dark energy and dark matter; a census of the solar system; the transient and variable optical sky; and the structure of the Milky Way \citep[][]{2019ApJ...873..111I}.
An automated \gls{Data Management System} processes tens of terabytes of images per night to produce science-ready data products within minutes of each observation for a global community of scientists.

This paper presents Rubin's third Data Preview, \gls{DP2} \citep[][]{10.71929/rubin/2570308}, and the first to be based entirely on LSSTCam data.
It marks the conclusion of the Rubin commissioning program and the transition to full survey operations.
\gls{DP2} is a reprocessing of all science-grade exposures acquired with LSSTCam during on-sky commissioning and the initial phase of early operations, prior to the formal start of the \gls{LSST}.
As such, it represents a natural boundary and provides a comprehensive dataset that captures the on-sky performance of the as-built system before LSST starts.

Two earlier Data Previews preceded \gls{DP2} as part of the Rubin Early Science Program \citep[][]{RTN-011}, established to provide the community with early access to data products and services needed to produce high-impact science during the period spanning commissioning through the first Data Release.
DP0 was based on simulated \gls{LSST}-like data \citep[][]{2021ApJS..253...31L} and served primarily to stand up early versions of the RSP and prepare the community for data access at scale.
DP1 \citep[][]{10.71929/rubin/2570308} delivered the first data products derived from on-sky observations using science-grade exposures acquired with LSSTComCam \citep[][]{2022SPIE12184E..0JS,2020SPIE11447E..0LS,2018SPIE10700E..3DH, 10.71929/rubin/2561361}; a 144-megapixel, single-raft instrument used during early commissioning prior to the installation of LSSTCam.
\gls{DP2} supersedes both in scope, fidelity, and scientific depth.

\gls{DP2} contains \nexposures science-grade exposures acquired over \svcampaignnnights nights between \svcampaignstartdate and \svcampaignenddate, covering a total sky area of \totalarea across \nfields fields.
The data products include raw and calibrated single-epoch images, coadded images, difference images, source detection catalogs, and a range of derived data products.
The release totals approximately \sizeinbytes\ and contains around \nobjects\ distinct astronomical objects.
Full release documentation is available at \url{https://dp2.lsst.io}.

gls{DP2}, like all Rubin Data Previews and Data Releases, is subject to a 2,yr proprietary period; during this time access is restricted to \gls{LSST} data rights holders \citep[][]{rdo-013} via the \gls{RSP} \citep[][]{LSE-319}, Rubin's cloud-based data science platform \citep[][]{LSE-319}.
After the proprietary period, \gls{DP2} will be made public, though access through Rubin Data Access Centers in the US and Chile will remain limited to data rights holders \citep[][]{rdo-013}.
Alternative public access mechanisms have not yet been finalized.

The structure of this paper is as follows.
Section~\ref{sec:commissioning} describes the observatory system and its status at the close of commissioning, the fields included in \gls{DP2}, and the observing strategy employed.
Section~\ref{sec:data_products} summarizes the data products in the release.
The processing pipelines are described in Section~\ref{sec:drp}, followed by data validation and performance characterization in Section~\ref{sec:performance}.
Section~\ref{sec:data_services} describes the RSP.
Section~\ref{sec:community_science} outlines Rubin's community support model.
A summary and outlook toward the start of formal \gls{LSST} Data Releases is given in Section~\ref{sec:summary}.

All magnitudes are quoted in the AB system \citep[][]{1983ApJ...266..713O} unless otherwise specified.